\documentclass[11pt]{article}
\usepackage{notes}
\usepackage{subfiles}

\begin{document}

% ---------- course title page ----------
\begin{titlepage}
\centering
\vspace*{3cm}
{\Huge\bfseries\color{coursedark} PHYS-E0510 -- Advanced Computational Methods in Physics D\par}
\vspace{0.8cm}
{\Large\color{courseaccent} Lecture Notes\par}
\vspace{0.3cm}
{\large Sessions 1--5\par}
\vspace{2cm}
{\large Jose Lado\par}
\vspace{0.5cm}
{\large Department of Applied Physics, School of Science\\Aalto University, Finland\par}
\vspace{0.5cm}
{\large 2026\par}
\vfill
{\small Each session's title box below names its companion Jupyter notebook and
prerequisites.\par}
\end{titlepage}

\section*{How this course is organized}
\addcontentsline{toc}{section}{\textbf{How this course is organized}}

Every topic in this course comes to you in four different forms, and they are meant to be used
together rather than as substitutes for one another.

\begin{itemize}[nosep]
  \item \textbf{Recordings.} Each session is recorded; this is the lecture itself, delivered
    live, with all the informal asides, questions, and derivations on the fly that never quite
    make it into a written document.
  \item \textbf{Slides.} The slide deck for each session is the visual backbone of the lecture:
    the sequence of definitions, figures, and equations in the order they were actually
    presented. Use it to navigate a recording or to jog your memory before an exam.
  \item \textbf{Lecture notes.} What you are reading now. The notes follow the same session
    structure as the recordings and slides, but write out the physics in full, connected
    prose, with every derivation carried through by hand rather than compressed onto a slide.
    Read them either instead of watching a session again, or alongside the slides to fill in the
    steps a bullet point necessarily skips.
  \item \textbf{Jupyter notebooks.} Every computational session has a companion notebook,
    referenced in the title box at the start of each session below and flagged again wherever a
    \emph{Computational companion} box appears in the margin. This is where the physics actually
    gets computed: nothing here is a canned figure, you are meant to run the cells yourself,
    change parameters, and watch what happens.
\end{itemize}

None of these four are complete on their own. The recordings and slides give you the narrative
and the intuition; the lecture notes give you the derivations you can check line by line; the
notebooks give you the numerics to build intuition none of the others can give you by
themselves. A good way to work through a session is to watch the recording once, read the
matching notes with the notebook open alongside it, and then go back and tinker with the
notebook until the results stop feeling like something that just happened on someone else's
screen.

\subsection*{The computational libraries}

All of the numerics in this course are delegated to two Python libraries written and maintained
by the instructor, rather than to code written from scratch inside the notebooks. Both need to be
installed separately before the notebooks will run; see the notebooks' first cell for the exact
environment variables they expect.

\begin{itemize}[nosep]
  \item \textbf{\texttt{pyqula}} (\url{https://github.com/joselado/pyqula}) handles
    single-particle and mean-field electronic structure: tight-binding Hamiltonians, band
    structures, and the self-consistent Hubbard/CDW calculations of Session 1. It is what
    \texttt{classical\_magnetism.ipynb} is built on.
  \item \textbf{\texttt{dmrgpy}} (\url{https://github.com/joselado/dmrgpy}) handles everything
    many-body: exact diagonalization of spin and fermionic Hamiltonians, and the matrix-product-state
    and DMRG calculations of the later sessions. It is what the quantum-magnetism,
    interacting-fermion, and MPS notebooks are built on.
\end{itemize}

Both libraries are used purely as tools here: you are not expected to read their source to follow
this course, only to call the functions the notebooks show you. If you do want to look under the
hood, the reference documentation collected for this course is a good place to start.

\clearpage
\tableofcontents
\clearpage

\phantomsection\addcontentsline{toc}{section}{\textbf{Session 1: Classical Magnetism and Mean-Field Theory}}
\subfile{classical_magnetism}

\clearpage

\phantomsection\addcontentsline{toc}{section}{\textbf{Session 2: Exact Diagonalization of the Heisenberg Spin Chain}}
\subfile{quantum_magnetism}

\clearpage

\phantomsection\addcontentsline{toc}{section}{\textbf{Session 3: Exact Diagonalization in the Fermionic Basis}}
\subfile{many_body_fermions}

\clearpage

\phantomsection\addcontentsline{toc}{section}{\textbf{Session 4: Matrix Product States for Spin Chains}}
\subfile{mps_spin_chains}

\clearpage

\phantomsection\addcontentsline{toc}{section}{\textbf{Session 5: Matrix Product States for Fermionic Models}}
\subfile{mps_fermionic}

\end{document}
